\section{Ética, normativa e impacto}
\label{sec:etica}

\textbf{¿Sufre el insecto?} No hay consenso, y eso obliga a la precaución. Una revisión de más de 300 estudios
concluye que las cucarachas adultas cumplen seis de los ocho criterios usados para evaluar
dolor~\cite{gibbons2022,crump2023}; otros autores lo consideran improbable~\cite{adamo2016}. Ningún artículo de
insectos cíborg mide estrés durante la estimulación ni reporta el destino de los animales, y esa es la crítica más
dura al campo~\cite{dodd2014,mestre2024}.

\textbf{Norma en México.} La NOM-062-ZOO-1999 no menciona invertebrados~\cite{nom062}; aun así conviene someter el
protocolo a un comité de forma voluntaria. La Ley General de Vida Silvestre sí aplica: la especie es exótica
(artículo~3), su manejo solo puede hacerse en confinamiento (artículo~27) y hacerlo fuera es una infracción
(artículo~122)~\cite{lgvs}. \textbf{Un despliegue real en escombros requeriría una autorización que hoy no existe.}
Las pilas de litio son residuos de manejo especial: cada mochila debe recuperarse.

\textbf{Protocolo de trato.}
\begin{enumerate}
  \item \textbf{Reemplazar, reducir y refinar.} Lo que pueda probarse en simulación o con el circuito en el lazo se
        prueba así; el número de insectos se fija antes; sin cirugía, con mochila adherida y estímulo de luz.
  \item Mochila de 3~g como máximo, que se retira al terminar; se verifica que el insecto camina y se endereza.
  \item No más de 12 estímulos acumulados; sesiones de una hora y descanso de 2~h, como en~\cite{gracidas2026}.
  \item Solo machos adultos, doble contención y conteo de entrada y salida. Nunca se liberan.
  \item El insecto que no se endereza, deja de comer o muestra una lesión sale del estudio; la congelación no es un
        método de eutanasia válido~\cite{tucker2023}. El reporte declara bajas y destino final~\cite{fischer2024}.
\end{enumerate}

\begin{table}[!htb]
\centering\footnotesize
\caption{Impacto en los contextos del resultado de aprendizaje 4 de ABET y el atributo 5 de
CACEI~\cite{abet,cacei}.}
\label{tab:impacto}
\rowcolors{2}{SoftGray}{white}
\begin{tabularx}{\textwidth}{@{}>{\bfseries\raggedright\arraybackslash}p{1.9cm} >{\raggedright\arraybackslash}X >{\raggedright\arraybackslash}X@{}}
\toprule
\rowcolor{TableHead} Contexto & Riesgo & Cómo se atiende \\
\midrule
Global & Doble uso: la misma mochila puede vigilar & Fin humanitario; imagen de baja resolución; sin audio ni grabación \\
Económico & Costo que oculta cría, personal y vida útil corta & Reportar el costo por hora de búsqueda \\
Ambiental & Escape de una especie exótica y pilas abandonadas & Solo machos, doble contención, mochilas numeradas \\
Social & Rechazo de la víctima y privacidad en viviendas & Baliza que identifique al dispositivo; 64 píxeles no reconocen a nadie \\
\bottomrule
\end{tabularx}
\end{table}

El proyecto responde a la meta 11.5 de los ODS y a la prioridad~4 del Marco de Sendai~\cite{ods,sendai}. El Código
de Ética del IEEE obliga a sustentar las afirmaciones en los datos y a revelar lo que pueda poner en peligro al
público~\cite{ieeecode}; por eso este documento separa lo medido, lo simulado y lo supuesto.
